% SHARD-ID: AQM-30-SPEC-FQ-XY-ARITHMETIC-FIELD
% SHARD-TITLE: The Arithmetic Field on \(\operatorname{Spec}\mathbb F_q[x,y]\)
% SHARD-SUMMARY: Computes the point types of the affine plane over \(\mathbb F_q\): generic, curve-generic, and closed points.
% SHARD-SUMMARY: Builds the closed-point Weyl-Heisenberg net, local algebras, and state presheaf for the affine plane.
% SHARD-SUMMARY: Separates rational-point truncations from the full closed-point quasi-local algebra.
% SHARD-KEYWORDS: affine plane, finite field, closed point, maximal ideal, curve point, Weyl-Heisenberg, rational truncation

\section{The Arithmetic Field on \texorpdfstring{\(\operatorname{Spec}\mathbb F_q[x,y]\)}{Spec Fq[x,y]}}
\label{sec:spec-fq-xy-arithmetic-field}

\subsection{Status and Source Anchors}

This shard instantiates AQM-28 for \(R_2=k[x,y]\), \(k=\mathbb F_q\).  The
Stacks Project source describes \(\Spec k[x,y]\) in
\path{references/algebraic_geometry/stacks_project_algebra.tex}, lines
4889--4927: besides \((0)\), height-one primes are generated by irreducible
polynomials, and non-principal primes are maximal of the displayed quotient
form.  The finite residue field statement for maximal ideals is the
Nullstellensatz source, lines 6520--6665.  Weyl and state conventions are
those of AQM-28.

\subsection{Lamport Dependency Skeleton}

\[
\begin{array}{rcl}
D1&:&R_2=k[x,y],\quad X_2=\Spec R_2.\\
L1&:&X_2\text{ has }(0),\text{ height-one points }(f),\text{ and
      closed maximal ideals }\mathfrak m.\\
L2&:&\kappa(0)=k(x,y),\
      \kappa((f))=\operatorname{Frac}(k[x,y]/(f)),\
      \kappa(\mathfrak m)/k\text{ finite}.\\
D2&:&C_2(U)=U\cap\operatorname{Max}R_2,\quad
      E_2(U)=\bigoplus_{\mathfrak m\in C_2(U)}\kappa(\mathfrak m)^2.\\
T1&:&\mathcal A_2(U)=
      \varinjlim_{S\Subset C_2(U)}
      \bigotimes_{\mathfrak m\in S}M_{q^{d(\mathfrak m)}}(\mathbb C).\\
T2&:&\text{The rational-point truncation is a finite }q^2\text{-site
      subtheory, not the full Spec theory.}\\
T3&:&\operatorname{State}(\mathcal A_2(U))\text{ is an inverse limit over
      finite closed supports.}\\
T4&:&\text{Coordinate profiles such as }x,y\text{ are formal
      infinite-support fields.}
\end{array}
\]

\subsection{Point Types L1--L2}

Let \(X_2=\Spec k[x,y]\).

\paragraph{Lemma \(L1\).}
The points of \(X_2\) have the following forms.
\[
\begin{array}{ll}
\text{generic point:} & (0),\\
\text{curve-generic points:} & (f),\quad f\in k[x,y]\text{ irreducible},\\
\text{closed points:} & \mathfrak m\subset k[x,y]\text{ maximal}.
\end{array}
\]

\emph{Proof.}
The Stacks example cited above records that \((0)\) is prime, that a principal
ideal generated by an irreducible polynomial is prime because \(k[x,y]\) is a
UFD, and that the remaining non-principal primes are maximal after reducing
to a prime of a one-variable polynomial ring over a finite algebraic
extension of \(k\).  This is exactly the displayed trichotomy.

\paragraph{Lemma \(L2\).}
The residue fields are
\[
  \kappa((0))=k(x,y),
\]
\[
  \kappa((f))=\operatorname{Frac}(k[x,y]/(f))
  \quad\text{for irreducible }f,
\]
and, for a closed point \(\mathfrak m\),
\[
  \kappa(\mathfrak m)=k[x,y]/\mathfrak m
  \quad\text{is a finite extension of }k.
\]
If \(d(\mathfrak m)=[\kappa(\mathfrak m):k]\), then
\(|\kappa(\mathfrak m)|=q^{d(\mathfrak m)}\).

\emph{Proof.}
The first two formulas are the definition of residue field as the fraction
field of \(R/\mathfrak p\).  The maximal-ideal formula follows because
localizing at a maximal ideal does not change the quotient residue field.
The Nullstellensatz source gives that this residue field is finite over
\(k\), and the cardinality follows from its \(k\)-vector-space dimension.

The generic point and curve-generic points are not finite Weyl sites in the
current convention.  The generic residue \(k(x,y)\) is infinite.  If
\(f\in k[x]\) is irreducible, then
\[
  k[x,y]/(f)\cong (k[x]/(f))[y],
\]
so \(\kappa((f))\cong (k[x]/(f))(y)\), also infinite.  If \(f\) has positive
degree in \(y\), then \(k[x]\hookrightarrow k[x,y]/(f)\), since a nonzero
polynomial in \(k[x]\) cannot be divisible by such an \(f\).  Hence the
quotient and its fraction field are infinite.  After swapping \(x\) and \(y\)
if necessary, this covers every irreducible \(f\).

\subsection{Closed Points as Finite Qudit Sites}

A rational point \((a,b)\in k^2\) gives the maximal ideal
\[
  \mathfrak m_{a,b}=(x-a,y-b),
\]
with residue field \(k\).  Hence it is one \(q\)-level site.

Closed points of larger degree are equally concrete.  Let \(\pi(x)\in k[x]\)
be monic irreducible of degree \(d\), let \(K=k[x]/(\pi)\), and choose
\(\beta\in K\) represented by \(\widetilde\beta(x)\in k[x]\).  Then
\[
  \mathfrak m_{\pi,\beta}=(\pi(x),\,y-\widetilde\beta(x))
\]
is maximal, with residue field \(K\).  It contributes one \(q^d\)-level
finite-field qudit.

\emph{Derivation.}
The quotient by \((\pi(x),y-\widetilde\beta(x))\) sends \(x\) to the class of
\(x\) in \(K\) and \(y\) to \(\beta\), giving a quotient isomorphic to \(K\).
Since \(K\) is a field, the ideal is maximal.

\subsection{Zariski Opens and Closed-Point Supports}

For \(h\in k[x,y]\),
\[
  D(h)=\{\mathfrak p:h\notin\mathfrak p\}.
\]
The closed-point support used by the quantum net is
\[
  C_2(D(h))=
  \{\mathfrak m\in\operatorname{Max}R_2:h\notin\mathfrak m\}.
\]
Equivalently, if \(\bar x,\bar y\) are the residue classes in
\(\kappa(\mathfrak m)\), then
\[
  \mathfrak m\in C_2(D(h))
  \quad\Longleftrightarrow\quad
  h(\bar x,\bar y)\ne0.
\]

Examples:
\[
  C_2(D(x))=\{\mathfrak m:x\notin\mathfrak m\},
\]
so \(D(x)\) removes the closed points on the vertical line \(x=0\).  The open
\[
  D(x)\cup D(y)=X_2\setminus V(x,y)
\]
removes only the rational origin \((x,y)\) from the closed-point support.

\subsection{Weyl-Heisenberg Net T1}

For an open \(U\subset X_2\), define
\[
  E_2(U)=\bigoplus_{\mathfrak m\in C_2(U)}\kappa(\mathfrak m)^2.
\]
For a finite support \(S\Subset C_2(U)\),
\[
  \mathcal A(S)=
  \bigotimes_{\mathfrak m\in S}
  \operatorname{End}_{\mathbb C}\ell^2(\kappa(\mathfrak m))
  \cong
  \bigotimes_{\mathfrak m\in S}M_{q^{d(\mathfrak m)}}(\mathbb C).
\]
The open-local algebra is
\[
  \mathcal A_2(U)=
  \varinjlim_{S\Subset C_2(U)}\mathcal A(S)
  =
  \operatorname{span}_{\mathbb C}\{W_U(e):e\in E_2(U)\}.
\]

\paragraph{Theorem \(T1\).}
\(U\mapsto\mathcal A_2(U)\) is the affine-plane open-local
Weyl-Heisenberg net.

\emph{Proof.}
This is AQM-28 with \(R=R_2\).  Inclusion of opens gives inclusion of closed
supports, hence tensor-by-identity inclusions of finite matrix algebras and an
inclusion of their filtered colimits.  The Weyl commutator at a closed point
\(\mathfrak m\) is
\[
  Z_{\mathfrak m}(b)X_{\mathfrak m}(a)
  =
  \psi_{\mathfrak m}(ab)X_{\mathfrak m}(a)Z_{\mathfrak m}(b),
\]
and different closed points commute.

\subsection{Rational-Point Truncation T2}

Let
\[
  S_{\rm rat}=\{\mathfrak m_{a,b}:a,b\in k\}.
\]
This finite support has \(q^2\) closed points, each with residue field \(k\).
The associated finite subtheory is
\[
  E(S_{\rm rat})=\bigoplus_{(a,b)\in k^2}k^2,
\]
\[
  \mathcal H(S_{\rm rat})=\bigotimes_{(a,b)\in k^2}\ell^2(k),
  \qquad
  \mathcal A(S_{\rm rat})\cong
  \bigotimes_{(a,b)\in k^2}M_q(\mathbb C).
\]
This is the finite \(k\)-rational-point truncation, matching the finite-set
construction on \(k^2\).  It is not the full \(\Spec k[x,y]\) theory, because
the full closed-point set also contains all non-rational closed points of
degrees \(d>1\).

\subsection{States T3}

For every open \(U\),
\[
  \operatorname{State}(\mathcal A_2(U))
  \cong
  \varprojlim_{S\Subset C_2(U)}
  \operatorname{State}(\mathcal A(S)).
\]
The transition maps are finite partial traces.

The Zariski state presheaf is not separated.  Let
\[
  x_0=(x,y),\qquad x_1=(x-1,y),
\]
and let \(U=X_2\setminus\{x_0\}\), \(V=X_2\setminus\{x_1\}\).  These are open
and cover \(X_2\).  On the two rational \(q\)-level sites \(x_0,x_1\), use
the same maximally entangled density matrix and product density matrix as in
AQM-29, tensoring with the product tracial state on every other closed point.
The two global states agree after restriction to \(U\) and to \(V\), but they
differ on the two-site algebra.  Thus local marginals on a Zariski cover do
not determine global correlations.

\subsection{Coordinate Field Profiles T4}

The coordinate pair \((x,0)\) defines the closed-point profile
\[
  \mathfrak m\longmapsto (x\bmod\mathfrak m,0).
\]
It is nonzero at every closed point not lying on the line \(x=0\).  Hence it
has infinite support: the line \(x=1\) has the same closed points as
\(\Spec k[y]\), and AQM-29 proves there are infinitely many of them.  Likewise
\((0,y)\) has infinite support, using the line \(y=1\).  Their
finite-volume restrictions define Weyl operators on finite closed supports,
but the full coordinate profiles are not elements of
\(\mathcal A_2(X_2)\).

The curve-generic point \((x)\) is also not a replacement for the missing
operator: its residue field is \(k(y)\), not a finite qudit field.  It records
the generic point of the divisor \(x=0\), while the finite Weyl theory records
finite closed-point excitations on or off that divisor.

\subsection{Conclusion}

\(\Spec k[x,y]\) has three geometric point types.  The finite arithmetic
quantum field constructed here uses only the closed points as finite qudit
sites, with one \(q^{d(\mathfrak m)}\)-level qudit at a closed point of degree
\(d(\mathfrak m)\).  Rational points give a finite \(q^2\)-site truncation,
but the full closed-point quasi-local algebra is the filtered tensor product
over all closed points of all degrees.  Generic and curve-generic points
remain essential to the Zariski topology and to divisor geometry, but they
are not finite Weyl sites under the current convention.
